% Shared abstract (included by main.tex, ieee.tex, springer.tex)
Many systems chain an AI model's typed decisions into a workflow and use
the model's reported confidence to decide which cases a human should
check. We test how well that works, using Jev, TypeSafe's typed decision
model and its three primitives (a categorical \choice{}, a binary
\noul{} and an ordinal \score{}), on a bug-triage workflow. We ask
whether these signals identify erroneous chained decisions, and whether
they identify the errors worth human review. We study 442 Mozilla and
Linux kernel bugs, then run pre-registered
tests on 500 new bugs, four further models and an independent Kubernetes
workflow. Failure is disagreement with a target action we construct from
reference labels. A frozen decision-feature estimator ranks failures
better than final-step confidence (AUROC 0.76 vs.\ 0.55 on the new bugs),
and this holds for all five models. Most of its signal is simply how often
the chosen action has been wrong before. The Kubernetes replication,
on a chronological split, confirmed this (AUROC 0.78 vs.\ 0.63), as did
three of four further models, including the only other one competent at
every step. These error rates do not transfer across label mixes: fitted
on a quota-sampled development set, they can rank natural-mix errors below
chance. Re-weighting to the deployment mix repairs the ranking, but the
label-free estimate of that mix is biased and made ranking worse on
Kubernetes. Better failure ranking did not reduce cost. Under the
pre-registered costs, no 10\% review policy beat no review. An oracle
shows that the savings sit in wrongly accepted cases, which the tested
signals could not pick out, and whether review pays off depends on how
costly a wrong accept is compared with a review.
